% =====================================================================
% DIAPOSITIVAS — Sustentación · Entrega 1 · Anteproyecto (10 minutos)
% Consultoría Estadística USTA 2026-II · Beamer
% =====================================================================
\documentclass[10pt,aspectratio=169]{beamer}

\usepackage[utf8]{inputenc}
\usepackage[T1]{fontenc}
\usepackage[spanish]{babel}
\usepackage{lmodern}
\usepackage{booktabs}
\usepackage{tabularx}

\usetheme{Madrid}
\usecolortheme{default}
\definecolor{usta}{RGB}{30,55,128}
\definecolor{accent}{RGB}{183,135,30}
\setbeamercolor{structure}{fg=usta}
\setbeamercolor{frametitle}{bg=usta,fg=white}
\setbeamercolor{title}{fg=usta}
\setbeamercolor{block title}{bg=usta,fg=white}
\setbeamertemplate{navigation symbols}{}
\setbeamertemplate{footline}{%
  \hfill\insertframenumber/\inserttotalframenumber\hspace{2mm}\vspace{1mm}}

\title[Auditoría del padrón MinCiencias]{Auditoría estadística reproducible del sistema de reconocimiento de investigadores de MinCiencias (2013--2021) con datos abiertos}
\subtitle{Entrega 1 — Anteproyecto · Consultoría Estadística · USTA 2026-II}
\author[Chaparro · Muñoz · Triana]{Kevin Leonardo Chaparro Reyes \and Valentina Muñoz Palma \and Paula Margarita Triana Ancinez}
\institute{Universidad Santo Tomás}
\date{\today}

\begin{document}

\begin{frame}[plain]
\titlepage
\end{frame}

\begin{frame}{El problema: ¿qué audita un instrumento público de medición?}
\begin{itemize}
  \item MinCiencias reconoce por \textbf{convocatorias} (2013--2021) a los investigadores del país y los clasifica: Junior, Asociado, Senior, Emérito.
  \item Ese padrón es la \textbf{principal fuente administrativa de capital humano científico} y alimenta decisiones de política, financiación y carrera.
  \item \textbf{Problema doble:} (1) nadie audita la \emph{calidad y comparabilidad} del registro; (2) la \emph{inequidad} estructural no está cuantificada con rigor en una sola fuente.
\end{itemize}
\vspace{2mm}
\textbf{A quién le importa:} MinCiencias y el SNCTI · los investigadores · los tomadores de decisión de política de CTI.
\end{frame}

\begin{frame}{Evidencia propia de que el problema existe}
\begin{itemize}
  \item Padrón público declara \textbf{77\,237 registros / 6 convocatorias} (30\,086 únicos)… pero el archivo consolidado verificado tiene \textbf{50\,891 filas y 3 convocatorias} (781/2017, 833/2018 y 894/2021) con \textbf{26\,662 personas únicas}.
  \item \textbf{22 registros con edad $>$ 100 años} (máx. 956) y errores de codificación.
  \item Concentración territorial: \textbf{Bogotá + Antioquia = 51\,\%}; Bogotá retiene 92\,\% y absorbe 1\,000+ investigadores.
  \item Brecha de género por área: \textbf{24\,\% mujeres en Ingeniería} vs.\ \textbf{48\,\% en ciencias médicas}.
  \item Vs.\ censo DANE 2018: afro \textbf{3$\times$}, indígena \textbf{7,9$\times$}, discapacidad \textbf{8$\times$} subrepresentados.
  \item Solo \textbf{36,9\,\%} de los autores de la producción declarada son investigadores reconocidos.
\end{itemize}
\end{frame}

\begin{frame}{La pregunta que vamos a responder}
\begin{block}{Pregunta estadística respondible}
\textbf{¿El padrón de investigadores reconocidos de MinCiencias (2013--2021) es fiable, comparable entre convocatorias y equitativo, y qué mejoras metodológicas y de gobernanza de datos lo convertirían en la base sólida de un observatorio de capital humano científico?}
\end{block}
\vspace{3mm}
\textbf{Si nadie hace nada:} decisiones de política y de carrera seguirán apoyándose en un instrumento cuya calidad no está auditada de forma reproducible.
\end{frame}

\begin{frame}{Estado del arte: lo que ya se sabe}
\begin{columns}[T]
\column{0.58\linewidth}
\textbf{Enfoques (mundo):}
\begin{itemize}
  \item \emph{Science of science} y métricas responsables (Fortunato 2018; Leiden 2015; CoARA 2022)
  \item Calidad de datos y reproducibilidad (Wang \& Strong 1996; FAIR 2016)
  \item Carreras longitudinales: Markov, supervivencia, movilidad (Sinatra 2016; Deville 2014)
  \item Desigualdad (HHI/Gini/Theil), género (Larivière 2013), redes y jerarquías (Clauset 2015)
\end{itemize}
\column{0.42\linewidth}
\textbf{En Colombia / región:}
\begin{itemize}
  \item Lattes (Brasil) y ScienTI-CVLAC/Publindex (Colombia)
  \item Críticas al modelo de categorización del SNCTI
  \item CONPES 4069 (ciencia abierta)
\end{itemize}
\vspace{1mm}
\textbf{Contraste:} cadenas de Markov (homogeneidad) vs.\ modelos de supervivencia (tiempo al evento).
\end{columns}
\vspace{2mm}
\emph{Vacío:} falta una \textbf{auditoría estadística reproducible} del padrón que combine calidad + longitudinal + equidad.
\end{frame}

\begin{frame}{Objetivos}
\textbf{General:} cuantificar y auditar de forma reproducible la fiabilidad, comparabilidad y equidad del padrón (2013--2021) y entregar un marco de indicadores y un tablero de auditoría.
\begin{enumerate}
  \item \textbf{Cuantificar la calidad} del registro (completitud, unicidad, validez, consistencia) por convocatoria.
  \item \textbf{Estimar la dinámica longitudinal}: retención, transición de categoría y tiempo hasta la salida (Markov + supervivencia).
  \item \textbf{Cuantificar la equidad}: concentración territorial (HHI/Gini/Theil) y representación de género, etnia, discapacidad y víctimas vs.\ DANE 2018.
\end{enumerate}
\emph{Cada objetivo declara con qué evidencia se dará por cumplido; la secuencia es visible: registro limpio $\rightarrow$ panel fiable $\rightarrow$ lectura de equidad.}
\end{frame}

\begin{frame}{Datos y método preliminar}
\begin{columns}[T]
\column{0.55\linewidth}
\textbf{Datos abiertos verificados (datos.gov.co):}
\begin{itemize}
  \item Padrón \texttt{bqtm-4y2h}: declara 6 convocatorias (2013--2021); verificadas 3 en el artefacto (781/2017, 833/2018 y 894/2021)
  \item Producción \texttt{33dq-ab5a}: 3,2 M filas (cruce por \texttt{id\_persona})
  \item Censo DANE 2018 (referencia de equidad)
\end{itemize}
\column{0.45\linewidth}
\textbf{Ruta en 3 etapas (Python + DuckDB, reproducible):}
\begin{enumerate}
  \item Auditoría de calidad (contratos de esquema)
  \item Panel longitudinal: Markov + supervivencia (KM/Cox)
  \item Equidad: HHI/Gini/Theil + razones de representación con IC
\end{enumerate}
\end{columns}
\vspace{2mm}
\textbf{Alternativa descartada:} análisis de producción individual ($h$); requiere ~1,2\,GB y queda como extensión. \textbf{Plan si el dato no llega:} OE1--OE3 no dependen de ese dataset; redescarga por API.
\end{frame}

\begin{frame}{Alcance y riesgos}
\textbf{Qué no haremos (y por qué):}
\begin{itemize}
  \item No validaremos contra registros internos de MinCiencias (no accesibles): se audita la fuente pública declarada.
  \item No haremos modelos predictivos individuales ni análisis causal (fuera del alcance semestral).
  \item No profundizaremos en producción individual ($h$) en esta fase.
\end{itemize}
\textbf{Riesgos con contingencia:}
\begin{enumerate}
  \item Discrepancia 77\,237 vs.\ 50\,891 $\rightarrow$ es un \emph{hallazgo del objetivo}, no un bloqueo.
  \item Producción no descargable $\rightarrow$ OE1--OE3 no dependen de ella.
  \item Tiempo ajustado $\rightarrow$ priorización por hitos + 2 semanas de holgura.
\end{enumerate}
\end{frame}

\begin{frame}{Por qué cabe en el semestre}
\textbf{Cronograma por hitos (12 semanas, con holgura):}
\begin{center}
\small
\begin{tabular}{ll}
\toprule
\textbf{Semanas} & \textbf{Hito} \\
\midrule
1--2 & Datos y contratos de esquema \\
3--4 & Auditoría de calidad (OE1) \\
5--6 & Panel y transiciones (OE2.a) \\
7--8 & Supervivencia (OE2.b) \\
9--10 & Equidad (OE3) \\
11 & Marco de indicadores + tablero \\
12 & Informe final y sustentación \\
\bottomrule
\end{tabular}
\end{center}
\vspace{2mm}
\textbf{Horas:} 3 integrantes $\times$ 6 h/semana $\times$ 12 semanas $\approx$ 216 h; carga estimada $\approx$ 180 h $\Rightarrow$ \textbf{~17\,\% de holgura}.
\end{frame}

\begin{frame}{Cierre}
\begin{itemize}
  \item El problema es real y \textbf{ya medido} con cifras propias verificables.
  \item La literatura (mundo y Colombia) deja un \textbf{vacío claro}: falta la auditoría reproducible del registro.
  \item Los datos son \textbf{abiertos, accesibles y ya abiertos por el equipo}.
  \item El plan es \textbf{realista} (hitos, holgura, horas).
\end{itemize}
\vspace{4mm}
\begin{block}{Aporte}
Un marco de indicadores y un tablero de auditoría que permitan a MinCiencias y a la comunidad saber \emph{qué tan confiable, comparable y equitativo} es el sistema de reconocimiento.
\end{block}
\vspace{3mm}
\begin{center}\large\textbf{¡Gracias!}\end{center}
\end{frame}

\end{document}
